A thin modular reactive mat is investigated as a means of reducing contaminant flux from a contaminated sediment source. The software combines one-dimensional reactive transport in a keratin core, geotextile resistance, independent tile-condition variables, two-dimensional coastal transport, synthetic observations, interval estimates and maintenance recommendations. This report documents all operating assumptions, the supplied keratin literature, six regenerated scenarios, three servicing policies and independent numerical checks.

In the six-year baseline, whole-hotspot attenuation is 61.32\% for Pb, 61.57\% for Hg and 61.19\% for Cu. The corresponding sorbed fractions of nominal capacity are 66.76\%, 0.83\% and 9.53\%. Physical resistance, chemical storage and effective area therefore require separate interpretation. Six-year cumulative whole-hotspot Pb emission is 642.3~kg without a mat, 62.74~kg with calendar servicing and 141.6~kg with evidence-informed servicing. These are synthetic source comparisons, not measured remediation yields.

The plot audit identified coordinate and area-weighting errors, off-by-one integration, out-of-range snapshot labels, an inconsistent barrier reference, and observation-clock and attribution defects. The corrected implementation was tested and all affected scenarios, figures and manuscript values regenerated. Remaining abrupt changes correspond to prescribed interventions, constitutive capacity locking or phase-dependent coastal statistics; numerical sensitivity is reported explicitly. The column and plume budgets close individually, but the area-mixed long-term source is not a fully coupled sediment--mat--water inventory.

The three supplied articles contain useful laboratory data but do not validate the finished mat in seawater. The default capacities of 1, 2.5 and 3~mg\,g$^{-1}$ for Pb, Hg and Cu remain operating assumptions. Laboratory mercury uptake by modified human hair and high batch removal by a keratin composite are documented with their different materials and experimental conditions. The resulting demonstrator supports hypothesis testing and experimental planning; field efficacy, ecological benefit and a reliable economic servicing advantage remain unestablished.
